\chapter{Algebraic and Analytic Background}
\label{chap:background}

This appendix collects several results from algebra and analysis that are used at various points throughout the report. The topics are diverse: field theory, complex analysis, elementary number theory, roots of unity, and symmetric polynomials, but each result is self-contained and included for ease of reference.

\section{Fields}

For any algebraic extension $E/\bbF_q$, the \emph{Galois group} $\gal(E/\bbF_q)$ is the group of all field automorphisms $\sigma: E \to E$ that fix every element of $\bbF_q$. For the extension $\bbF_{q^r}/\bbF_q$, this group is cyclic of order $r$, generated by the Frobenius automorphism $\phi: \alpha \mapsto \alpha^q$. Recall from the chapter on finite fields that, one of the key fact connecting automorphisms to irreducible polynomials and it's roots is the \IsoExtthm from the \autoref{subsec:splitting-field}: any isomorphism between two subfields of an algebraic closure extends to an automorphism of the whole closure, and in particular, any two roots of an irreducible polynomial over $\bbF_q$ are conjugates under the action of $\gal(\bbF_{q^r}/\bbF_q)$.

The following theorem makes this concrete for irreducible polynomials over $\bbF_q$. It describes precisely how an irreducible polynomial $h\in\bbF_q[X]$ of degree $d$ behaves when viewed over a larger field $\bbF_{q^k}$: it breaks into $\gcd(k,d)$ irreducible factors, all of the same degree, and the Galois group of $\bbF_{q^r}/\bbF_q$ (where $r=\gcd(k,d)$) acts transitively on these factors, exactly the statement of the Isomorphism Extension Theorem applied in this finite-field setting.

\begin{Theorem}{}{thm:irr-facts-into-irr-permutes}
  Let $K=\bbF_{q^k}$ and $E=\bbF_{q^r}$. Suppose $h(X)=X^d-a_1X^{d-1}+\cdots+(-1)^da_d$ be an irreducible polynomial in $\bbF_q[X]$. Then in $\bbF_{q^k}[X]$, $h$ splits into $r=\gcd(k,d)$ many irreducible polynomials of degree $\nicefrac{d}{r}$: $$h(X)=h_1(X)\cdots h_r(X)$$If we normalize $h_i(X)$ such that $h_i$ is monic then $h_i(X)\in\bbF_{q^r}[X]$ for all $i\in[r]$. Then for any $i,j\in[r]$ with $i\neq j$ there exists $\sg\in\gal(\bbF_{q^r}/\bbF_q)$ such that $\sg(h_i)=h_j$.
\end{Theorem}

\section{Meromorphic and Analytic Continuation}\label{sec:meromorphic-analytic}

The analytic continuation of the Riemann zeta function, and more generally of Dirichlet $L$-functions, requires the notion of a meromorphic function. We record the relevant definitions and a motivating example here.

\begin{Definition}{Zero, Pole, Order, Residue}{}
  \index{zero!of a function}\index{pole}\index{residue}%
  Assume $f$ has a convergent Laurent series expansion about $z_0$:
  $$f(z) = a_n(z-z_0)^n + a_{n+1}(z-z_0)^{n+1} + \cdots = \sum_{m=n}^{\infty} a_m(z-z_0)^m,$$
  with $a_n \neq 0$. If $n > 0$ we say $f$ has a \emph{zero of order $n$} at $z_0$; if $n < 0$ we say $f$ has a \emph{pole of order $-n$} at $z_0$. If $n = -1$ we say $f$ has a \emph{simple pole with residue $a_{-1}$}. We denote the order of $f$ at $z_0$ by $\ord_f(z_0) = n$.
\end{Definition}

\begin{Definition}{Meromorphic Function}{}
  \index{meromorphic function}\index{analytic function}%
  We say $f$ is \emph{meromorphic at $z_0$} if there exists a disk about $z_0$ of radius $r$ and an integer $n_0$ such that for all $z$ with $|z - z_0| < r$,
  $$f(z) = \sum_{n \geq n_0} a_n(z-z_0)^n.$$
  If $f$ is meromorphic at each point in a region, we say $f$ is \emph{meromorphic} on that region. If $n_0 \geq 0$, we say $f$ is \textbf{analytic}.
\end{Definition}

A canonical example illustrates the notion of continuation. Consider the geometric series $G(r) = \sum_{n=1}^{\infty} r^n$, which converges for $|r| < 1$ and equals $\frac{1}{1-r}$ there. Define $H(r) = \frac{1}{1-r}$; then $H$ is well-defined for all $r \neq 1$ and agrees with $G$ wherever $G$ converges. Since $H$ is defined on a strictly larger domain and has a simple pole at $r = 1$ (with residue $-1$), we call $H$ a \textbf{meromorphic continuation}\index{meromorphic continuation} of $G$. Had $H$ no poles, we would call it an \textbf{analytic continuation}\index{analytic continuation}. A function that is analytic at every point of $\bbC$, equivalently one whose Taylor expansion converges everywhere so that it has no poles anywhere, is said to be \textbf{entire}\index{entire function}.

\section{Fermat's Little Theorem}
Fermat's Little Theorem is one of the oldest results in elementary number theory, recorded in a 1640 letter from Pierre de Fermat to Fr\'{e}nicle de Bessy. It underlies essentially every computation we perform in the multiplicative group $\bbF_p^*$ -- from the very existence of the canonical additive character of $\bbF_q$ (where $q=p^r$), through the character congruence $\psi(l)\equiv l^f\bmod P$ used in the arithmetic proof of the Davenport--Hasse relation, down to the mod-$q$ reduction of Gauss sums appearing in the proof of the \QRthm. We record it here for completeness.
\begin{Theorem}{Fermat's Little Theorem}{thm:fermat-little}%
  \index{Fermat's little theorem}
  Let $p$ be a prime. Then for every integer $a$ we have $$a^p\equiv a\pmod p,$$ and in particular, if $p\nmid a$ then $a^{p-1}\equiv 1\pmod p$.
\end{Theorem}
\begin{proof}
  The non-zero residues modulo $p$ form a group $\bbF_p^*=\{1,2,\dots,p-1\}$ under multiplication of order $p-1$. By Lagrange's theorem every element of a finite group satisfies $g^{|G|}=e$, so $a^{p-1}\equiv 1\pmod p$ for every $a\in\bbF_p^*$, i.e. every $a$ with $p\nmid a$. Multiplying both sides by $a$ gives $a^p\equiv a\pmod p$ for such $a$; and the congruence $a^p\equiv a\pmod p$ is trivially true when $p\mid a$, so it holds for every integer $a$.
\end{proof}

\section{A Product Identity for Roots of Unity}

The following identity, used in the factorisation of lifted $L$-functions, expresses a product of linear factors twisted by roots of unity as a perfect power.

\begin{lemma}{Roots-of-Unity Product Identity}{lem:roots-unity-product}
  Let $n,m\in\bbN$ and set $d=\gcd(m,n)$. Then
  $$\prod_{t=1}^{n}\lt(1-e^{2\pi i\nicefrac{mt}{n}}\,X\rt)=\lt(1-X^{\nicefrac{n}{d}}\rt)^{d}.$$
\end{lemma}
\begin{proof}
  Let $\omega=e^{2\pi i/n}$ be a primitive $n$-th root of unity, so the factors are $1-\omega^{mt}X$ for $t=1,\dots,n$. Since $t=n$ gives $\omega^{mn}=1=\omega^0$, the product is the same as $\prod_{t=0}^{n-1}(1-\omega^{mt}X)$.

  Because $d=\gcd(m,n)$, the element $\omega^m$ is a primitive $(n/d)$-th root of unity. As $t$ ranges over $\{0,1,\dots,n-1\}$, the values $\omega^{mt}$ cycle through the $(n/d)$-th roots of unity exactly $d$ times each. Hence,
  $$\prod_{t=0}^{n-1}\lt(1-\omega^{mt}X\rt)=\lt(\prod_{k=0}^{n/d-1}\lt(1-(\omega^m)^k X\rt)\rt)^d.$$
  Since $\omega^m$ is a primitive $(n/d)$-th root of unity, the inner product equals $1-X^{n/d}$ by the standard factorisation of $1-Y^N = \prod_{k=0}^{N-1}(1-\zeta^k Y)$. Therefore, the whole product equals $(1-X^{n/d})^d$. \qed
\end{proof}

\section{Elementary Symmetric Polynomials}
Waring's Formula expresses each power-sum symmetric polynomial $s_k=x_1^k+x_2^k+\cdots+x_n^k$ as an explicit polynomial in the elementary symmetric polynomials $\sigma_1,\dots,\sigma_n$. It is the source of the closed-form reduction of lifted Kloosterman sums $K(\chi^{(s)};a,b)$ to the ground-field sum $K(\chi;a,b)$ in \lemref{lem:kloosterman-reduction}: applying Waring's formula to $\omega_1^s+\omega_2^s$ with $\sigma_1=\omega_1+\omega_2=-K(\chi;a,b)$ and $\sigma_2=\omega_1\omega_2=q$ produces the binomial-coefficient expansion. We state the general identity here for reference.
\begin{Theorem}{Waring's Formula}{thm:waring}%
  \index{Waring's formula}
  Let $x_1,\dots,x_n$ be indeterminates, let $$\sigma_j=\sum_{1\leq i_1<i_2<\cdots<i_j\leq n}x_{i_1}x_{i_2}\cdots x_{i_j}\qquad (j=1,\dots,n)$$ be the elementary symmetric polynomials in $x_1,\dots,x_n$, and let $$s_k=x_1^k+x_2^k+\cdots+x_n^k$$ be the $k^{\text{th}}$ power-sum. Then for every integer $k\geq 1$,
  $$s_k=\sum(-1)^{i_2+i_4+i_6+\cdots}\cdot \frac{(i_1+i_2+\cdots+i_n-1)!\cdot k}{i_1!\cdot i_2!\cdots i_n!}\cdot \sigma_1^{i_1}\sigma_2^{i_2}\cdots \sigma_n^{i_n},$$
  where the summation runs over all $n$-tuples $(i_1,\dots,i_n)$ of non-negative integers satisfying $i_1+2i_2+\cdots+ni_n=k$. Each coefficient of $\sigma_1^{i_1}\sigma_2^{i_2}\cdots \sigma_n^{i_n}$ is an integer.
\end{Theorem}

\begin{Theorem}{}{thm:elm-symmetry-spans}
  Suppose $f\in\bbF_q[X_1,\dots,X_n]$ is a symmetric polynomial. Let $\ovX=(X_1,\dots, X_n)$. Then there exists a polynomial $g\in \bbF_q[Y_1,\dots,Y_n]$ such that $$f(X_1,\dots, X_n)=g(\esym_1(\ovX),\dots,\esym_n(\ovX)).$$

  Moreover, \begin{enumerate}[label=(\roman*)]
    \item if $\deg_{X_i}(f)=d$ for all $i\in[n]$ then $\deg G=d$.
    \item if $\deg f=D$ then for every monomial $Y_1^{d_1}\cdots Y_n^{d_n}$  of $g$ with non-zero coefficient then $d_1+2\cdot d_2+\cdots +n\cdot d_n= D$.
  \end{enumerate}
\end{Theorem}
